% ====================================================================
%  徐森林、薛春华《数学分析》（第一册）
%  \S6.4  Riemann 积分的换元与分部积分
%  源：原书印刷页 386--398（PDF 物理页 392--404）
%  底本：PaddleOCR-VL-1.6，见 chapters/_ch06_4_raw.md
%  约定：全西文标点；行间公式用一对方括号；多行公式用 align 或 aligned；
%        跨册引用一律用 extref（写作「定理 x.y.z」原文）；数学模式内的中文必须包text
% ====================================================================
\section{Riemann 积分的换元与分部积分}\label{sec:6.4}

Newton-Leibniz 公式给了我们计算 Riemann 积分的简捷办法, 它的关键在于求出原函数. 换元 (变量代换) 与分部积分是求原函数或不定积分的两个重要方法. 用它们先求原函数再应用 Newton-Leibniz 公式计算 Riemann 积分, 但有时应用起来不免有些繁琐. 这里将要介绍的是不先求原函数, 而是套用换元与分部积分的方法与步骤直接计算 Riemann 积分.

\begin{theorem}[Riemann 积分的换元 (变量代换)]\label{thm:6.4.1}
设 $f(x)$ 在 $[a,b]$ 上连续, $x=\varphi(t)$ 在 $[\alpha,\beta]$ 上连续可导, 且 $\varphi([\alpha,\beta])\subset[a,b]$, $\varphi(a)=a$, $\varphi(\beta)=b$, 则
\[
\int_{a}^{b}f(x)\,\mathrm{d}x=\int_{a}^{\beta}f(\varphi(t))\varphi'(t)\,\mathrm{d}t=\int_{a}^{\beta}f\circ\varphi(t)\varphi'(t)\,\mathrm{d}t.
\]
\end{theorem}

\begin{proof}
因为 $f$ 在 $[a,b]$ 上连续, 根据微积分基本定理知, $f$ 在 $[a,b]$ 上有原函数 $F$, 即 $F'(x)=f(x)$, 所以,
\[
[F(\varphi(t))]'=F'(\varphi(t))\cdot\varphi'(t)=f(\varphi(t))\cdot\varphi'(t).
\]
再根据 Newton-Leibniz 公式得到
\begin{align*}
\int_{a}^{\beta}f(\varphi(t))\varphi'(t)\,\mathrm{d}t&=F(\varphi(t))\Big|_{\alpha}^{\beta}=F(\varphi(\beta))-F(\varphi(\alpha))\\
&=F(b)-F(a)=\int_{a}^{b}f(x)\,\mathrm{d}x.
\end{align*}
\end{proof}

如果在定理~\ref{thm:6.4.1} 中, 当 $t$ 从 $\alpha$ 变到 $\beta$ 时, $\varphi(t)$ 严格增地从 $a$ 变到 $b$, 则我们有不同特色的另一证明.

\begin{theoremp}[Riemann 积分的换元 (变量代换)]\label{thm:6.4.1p}
设 $f(x)$ 在 $[a,b]$ 上连续, $x=\varphi(t)$ 在 $[\alpha,\beta]$ 上连续可导, 且当 $t$ 从 $\alpha$ 到 $\beta$ 时, $\varphi(t)$ 严格递增地从 $a$ 变到 $b$, 则有
\[
\int_{a}^{b}f(x)\,\mathrm{d}x=\int_{a}^{\beta}f(\varphi(t))\varphi'(t)\,\mathrm{d}t=\int_{a}^{\beta}f\circ\varphi(t)\varphi'(t)\,\mathrm{d}t.
\]
\end{theoremp}

\begin{proof}
由于上式两边的被积函数都是连续的, 因此, Riemann 积分都存在.

作 $[\alpha,\beta]$ 上的分割 $W:\alpha=t_{0}<t_{1}<\cdots<t_{n}=\beta$. 令 $x_{i}=\varphi(t_{i})$, $i=0,1,\cdots,n$. 这就产生了 $[a,b]$ 的一个分割 $T:a=x_{0}<x_{1}<\cdots<x_{n}=b$. 由 Lagrange 中值定理, 我们有
\[
\Delta x_{i}=x_{i}-x_{i-1}=\varphi(t_{i})-\varphi(t_{i-1})=\varphi'(\tau_{i})(t_{i}-t_{i-1})=\varphi'(\tau_{i})\Delta t_{i},
\]
其中 $\tau_{i}\in(t_{i-1},t_{i})$, $i=1,2,\cdots,n$. 设 $K$ 为 $|\varphi'|$ 在 $[\alpha,\beta]$ 上的一个上界, 便有
\[
\Delta x_{i}\leqslant K\Delta t_{i},\quad i=1,2,\cdots,n.
\]
由此可知 $\|T\|\leqslant K\|W\|$. 所以, $\|W\|\to0$ 蕴涵着 $\|T\|\to0$.

令 $\xi_{i}=\varphi(\tau_{i})$, $i=1,2,\cdots,n$, 则易见 $x_{i-1}=\varphi(t_{i-1})<\xi_{i}=\varphi(\tau_{i})<\varphi(t_{i})=x_{i}$, 即 $\xi_{i}\in(x_{i-1},x_{i})$, $i=1,2,\cdots,n$. 于是,
\begin{align*}
\int_{a}^{b}f(x)\,\mathrm{d}x&=\lim_{\|T\|\to0}\sum_{i=1}^{n}f(\xi_{i})\Delta x_{i}=\lim_{\|W\|\to0}\sum_{i=1}^{n}f(\varphi(\tau_{i}))\varphi'(\tau_{i})\Delta t_{i}\\
&=\int_{a}^{\beta}f\circ\varphi(t)\varphi'(t)\,\mathrm{d}t.
\end{align*}
\end{proof}

\begin{theorempp}\label{thm:6.4.1pp}
在定理~\ref{thm:6.4.1p} 中, 如果将 $f(x)$ 在 $[a,b]$ 上「连续」改为「Riemann 可积」, 则结论仍正确.
\end{theorempp}

\begin{proof}
$\forall\varepsilon>0$, 因为 $\varphi'(t)$ 连续, 故 $\exists\delta_{1}>0$, 当 $t',t''\in[\alpha,\beta]$, $|t'-t''|<\delta_{1}$ 时, $|\varphi'(t')-\varphi'(t'')|<\dfrac{\varepsilon}{3M+1}$, 其中 $M=\sup\limits_{x\in[a,b]}|f(x)|$. 又因为 $f$ 在 $[a,b]$ 上 Riemann 可积, 故 $\exists\delta_{2}>0$, 当 $[a,b]$ 的分割 $T$, $\|T\|<\delta_{2}$ 时, $\sum\limits_{i=1}^{n}\omega_{i}^{f}\Delta x_{i}<\dfrac{\varepsilon}{3}$, 且
\[
\left|\sum_{i=1}^{n}f(\xi_{i})\Delta x_{i}-\int_{a}^{b}f(x)\,\mathrm{d}x\right|<\frac{\varepsilon}{3},\quad\xi_{i}\in[x_{i-1},x_{i}].
\]
于是, 对 $[\alpha,\beta]$ 的任一分割 $W:\alpha=t_{0}<t_{1}<\cdots<t_{n}=\beta$, $\forall\tilde{\tau}_{i}\in[t_{i-1},t_{i}]$, 令 $x_{i}=\varphi(t_{i})$, $x_{i}-x_{i-1}=\varphi(t_{i})-\varphi(t_{i-1})=\varphi'(\tau_{i})(t_{i}-t_{i-1})$, 当 $\|W\|<\delta=\min\left\{\delta_{1},\dfrac{\delta_{2}}{K+1}\right\}$ 时, 有 $\|T\|\leqslant K\|W\|<K\cdot\dfrac{\delta_{2}}{K+1}<\delta_{2}$, 其中 $K$ 为 $|\varphi'|$ 在 $[\alpha,\beta]$ 上的一个上界.
\begin{align*}
&\left|\sum_{i=1}^{n}f(\varphi(\tilde{\tau}_{i}))\varphi'(\tilde{\tau}_{i})\Delta t_{i}-\int_{a}^{b}f(x)\,\mathrm{d}x\right|\\
\leqslant{}&\sum_{i=1}^{n}\left|f(\varphi(\tilde{\tau}_{i}))\right|\bullet\left|\varphi'(\tilde{\tau}_{i})-\varphi'(\tau_{i})\right|\Delta t_{i}\\
&+\sum_{i=1}^{n}\left|f(\varphi(\tilde{\tau}_{i}))-f(\varphi(\tau_{i}))\right|\varphi'(\tau_{i})\Delta t_{i}\\
&+\left|\sum_{i=1}^{n}f(\varphi(\tau_{i}))\varphi'(\tau_{i})\Delta t_{i}-\int_{a}^{b}f(x)\,\mathrm{d}x\right|\\
\leqslant{}&M\frac{\varepsilon}{3M+1}+\sum_{i=1}^{n}\omega_{i}^{f}\Delta x_{i}+\left|\sum_{i=1}^{n}f(\xi_{i})\Delta x_{i}-\int_{a}^{b}f(x)\,\mathrm{d}x\right|\\
<{}&\frac{\varepsilon}{3}+\frac{\varepsilon}{3}+\frac{\varepsilon}{3}=\varepsilon.
\end{align*}
由此可知
\[
\int_{a}^{\beta}f(\varphi(t))\varphi'(t)\,\mathrm{d}t=\lim_{\|W\|\to0}\sum_{i=1}^{n}f(\varphi(\tilde{\tau}_{i}))\varphi'(\tilde{\tau}_{i})\Delta t_{i}=\int_{a}^{b}f(x)\,\mathrm{d}x.
\]
\end{proof}

\begin{theorem}[Riemann 积分的分部积分]\label{thm:6.4.2}
设 $u(x)$, $v(x)$ 为 $[a,b]$ 上的连续可导函数, 则有 Riemann 积分的分部积分公式
\[
\int_{a}^{b}u(x)v'(x)\,\mathrm{d}x=u(x)v(x)\Big|_{a}^{b}-\int_{a}^{b}u'(x)v(x)\,\mathrm{d}x,
\]
即
\[
\int_{a}^{b}u(x)\,\mathrm{d}v(x)=u(x)v(x)\Big|_{a}^{b}-\int_{a}^{b}v(x)\,\mathrm{d}u(x).
\]
\end{theorem}

\begin{proof}[证法 1]
\begin{align*}
&\int_{a}^{b}u(x)v'(x)\,\mathrm{d}x\\
&\xlongequal{\text{Newton-Leibniz 公式}}\left[u(x)v(x)-\int u'(x)v(x)\,\mathrm{d}x\right]\Big|_{a}^{b}\\
&=u(x)v(x)\Big|_{a}^{b}-\left[\int u'(x)v(x)\,\mathrm{d}x\right]\Big|_{a}^{b}\\
&=u(x)v(x)\Big|_{a}^{b}-\int_{a}^{b}u'(x)v(x)\,\mathrm{d}x.
\end{align*}
\end{proof}

\begin{proof}[证法 2]
因为 $uv$ 为 $uv'+u'v$ 在 $[a,b]$ 上的一个原函数, 故
\begin{align*}
&\int_{a}^{b}u(x)v'(x)\,\mathrm{d}x+\int_{a}^{b}u'(x)v(x)\,\mathrm{d}x\\
&=\int_{a}^{b}[u(x)v'(x)+u'(x)v(x)]\,\mathrm{d}x\\
&\xlongequal{\text{Newton-Leibniz}}u(x)v(x)\Big|_{a}^{b},
\end{align*}
从而
\[
\int_{a}^{b}u(x)v'(x)\,\mathrm{d}x=u(x)v(x)\Big|_{a}^{b}-\int_{a}^{b}u'(x)v(x)\,\mathrm{d}x.
\]
\end{proof}

进而, 有下面的定理.

\begin{theorem}\label{thm:6.4.3}
设 $u(x)$, $v(x)$ 在 $[a,b]$ 上具有 $n+1$ 阶连续导数, 则有
\begin{align*}
\int_{a}^{b}u(x)v^{(n+1)}(x)\,\mathrm{d}x&=\left[u(x)v^{(n)}(x)-u'(x)v^{(n-1)}(x)+\cdots+(-1)^{n}u^{(n)}(x)v(x)\right]\Big|_{a}^{b}\\
&\quad+(-1)^{n+1}\int_{a}^{b}u^{(n+1)}(x)v(x)\,\mathrm{d}x,\quad n=0,1,2,\cdots.
\end{align*}
\end{theorem}

\begin{proof}
\begin{align*}
\int_{a}^{b}u(x)v^{(n+1)}(x)\,\mathrm{d}x&=u(x)v^{(n)}(x)\Big|_{a}^{b}-\int_{a}^{b}u'(x)v^{(n)}(x)\,\mathrm{d}x\\
&=u(x)v^{(n)}(x)\Big|_{a}^{b}-u'(x)v^{(n-1)}(x)\Big|_{a}^{b}+\int_{a}^{b}u''(x)v^{(n-1)}(x)\,\mathrm{d}x\\
&=\cdots\\
&=\left[u(x)v^{(n)}(x)-u'(x)v^{(n-1)}(x)+\cdots+(-1)^{n}u^{(n)}(x)v(x)\right]\Bigg|_{a}^{b}\\
&\quad+(-1)^{(n+1)}\int_{a}^{b}u^{(n+1)}(x)v(x)\,\mathrm{d}x.
\end{align*}
\end{proof}

\begin{example}\label{ex:6.4.1}
Taylor 公式的积分型余项.

设函数 $f$ 在点 $x_{0}$ 的某开邻域 $U(x_{0})$ 内有 $n+1$ 阶连续导数. 令 $x\in U(x_{0})$, $u(t)=(x-t)^{n}$, $v(t)=f(t)$, $t\in[x_{0},x]\ (\text{或}\ t\in[x,x_{0}])$. 根据定理~\ref{thm:6.4.3}, 有
\begin{align*}
&\int_{x_{0}}^{x}(x-t)^{n}f^{(n+1)}(t)\,\mathrm{d}t\\
&=\left[(x-t)^{n}f^{(n)}(t)+n(x-t)^{n-1}f^{(n-1)}(t)+\cdots+n!f(t)\right]\Bigg|_{x_{0}}^{x}+\int_{x_{0}}^{x}0\cdot f(t)\,\mathrm{d}t\\
&=n!f(x)-n!\left[f(x_{0})+f'(x_{0})(x-x_{0})+\cdots+\frac{f^{(n)}(x_{0})}{n!}(x-x_{0})^{n}\right]\\
&=n!R_{n}(x),
\end{align*}
由此求得
\[
R_{n}(x)=\frac{1}{n!}\int_{x_{0}}^{x}f^{(n+1)}(t)(x-t)^{n}\,\mathrm{d}t,
\]
并称它为 Taylor 公式的积分型余项.

由于 $f^{(n+1)}(t)$ 连续, $(x-t)^{n}$ 在 $[x_{0},x]$ (或 $[x,x_{0}]$) 上保持同号, 因此, 由推广的积分第一中值定理, 有
\begin{align*}
R_{n}(x)&=\frac{1}{n!}\int_{x_{0}}^{x}f^{(n+1)}(t)(x-t)^{n}\,\mathrm{d}t\\
&=\frac{1}{n!}f^{(n+1)}(\xi)\int_{x_{0}}^{x}(x-t)^{n}\,\mathrm{d}t\\
&=\frac{1}{n!}f^{(n+1)}(\xi)\frac{-(x-t)^{n+1}}{n+1}\Bigg|_{x_{0}}^{x}\\
&=\frac{f^{(n+1)}(\xi)}{(n+1)!}(x-x_{0})^{n+1},\quad\xi\in(x_{0},x)\ (\text{或}\ (x,x_{0})),
\end{align*}
其中 $\xi=x_{0}+\theta(x-x_{0})$, $0<\theta<1$. 这就是熟悉的 Taylor 公式的 Lagrange 型余项.

如果直接应用积分第一中值定理, 就得到 Taylor 公式的 Cauchy 余项
\begin{align*}
R_{n}(x)&=\frac{1}{n!}\int_{x_{0}}^{x}f^{(n+1)}(t)(x-t)^{n}\,\mathrm{d}t\\
&=\frac{1}{n!}f^{(n+1)}(\xi)(x-\xi)^{n}\int_{x_{0}}^{x}\mathrm{d}t\\
&=\frac{1}{n!}f^{(n+1)}(\xi)(x-\xi)^{n}(x-x_{0})\\
&=\frac{1}{n!}f^{(n+1)}(x+\theta(x-x_{0}))(1-\theta)^{n}(x-x_{0})^{n+1},
\end{align*}
其中 $\xi=x_{0}+\theta(x-x_{0})$, $x-\xi=(x-x_{0})-\theta(x-x_{0})=(1-\theta)(x-x_{0})$.
\end{example}

\begin{example}\label{ex:6.4.2}
求下列 Riemann 积分:
\[
(1)\ \int_{0}^{1}\sqrt{1-x^{2}}\,\mathrm{d}x;\qquad(2)\ \int_{0}^{\frac{\pi}{2}}\sin t\cos^{2}t\,\mathrm{d}t;\qquad(3)\ \int_{1}^{\mathrm{e}}x^{2}\ln x\,\mathrm{d}x.
\]
\end{example}

\begin{solution}
(1)
\begin{align*}
\int_{0}^{1}\sqrt{1-x^{2}}\,\mathrm{d}x&\xlongequal{x=\sin t}\int_{0}^{\frac{\pi}{2}}\sqrt{1-\sin^{2}t}\cos t\,\mathrm{d}t\\
&=\int_{0}^{\frac{\pi}{2}}\cos^{2}t\,\mathrm{d}t=\frac{1}{2}\int_{0}^{\frac{\pi}{2}}(1+\cos2t)\,\mathrm{d}t\\
&=\frac{1}{2}\left(t+\frac{1}{2}\sin2t\right)\Bigg|_{0}^{\frac{\pi}{2}}=\frac{\pi}{4}.
\end{align*}
或者由下面的例~\ref{ex:6.4.5} 的结论, 得
\[
\int_{0}^{1}\sqrt{1-x^{2}}\,\mathrm{d}x=\int_{0}^{\frac{\pi}{2}}\cos^{2}t\,\mathrm{d}t=\frac{1}{2}\cdot\frac{\pi}{2}=\frac{\pi}{4}.
\]
(2)
\[
\int_{0}^{\frac{\pi}{2}}\sin t\cos^{2}t\,\mathrm{d}t\xlongequal{x=\cos t}\int_{1}^{0}x^{2}(-\mathrm{d}x)=-\frac{x^{3}}{3}\Bigg|_{1}^{0}=\frac{1}{3}.
\]
或者
\begin{align*}
\int_{0}^{\frac{\pi}{2}}\sin t\cos^{2}t\,\mathrm{d}t&=\int_{0}^{\frac{\pi}{2}}\sin t(1-\sin^{2}t)\,\mathrm{d}t\\
&=\int_{0}^{\frac{\pi}{2}}\sin t\,\mathrm{d}t-\int_{0}^{\frac{\pi}{2}}\sin^{3}t\,\mathrm{d}t=-\cos t\Bigg|_{0}^{\frac{\pi}{2}}-\frac{2!!}{3!!}=1-\frac{2}{3}=\frac{1}{3}.
\end{align*}
(3)
\begin{align*}
\int_{1}^{\mathrm{e}}x^{2}\ln x\,\mathrm{d}x&\xlongequal{\text{升幂}}\int_{1}^{\mathrm{e}}\ln x\,\mathrm{d}\left(\frac{x^{3}}{3}\right)\\
&\xlongequal{\text{分部积分}}\frac{1}{3}\left(x^{3}\ln x\Bigg|_{1}^{\mathrm{e}}-\int_{1}^{\mathrm{e}}x^{3}\cdot\frac{1}{x}\,\mathrm{d}x\right)\\
&=\frac{1}{3}\left(\mathrm{e}^{3}-\frac{x^{3}}{3}\Bigg|_{1}^{\mathrm{e}}\right)=\frac{1}{3}\left(\mathrm{e}^{3}-\frac{\mathrm{e}^{3}}{3}+\frac{1}{3}\right)=\frac{1}{9}(2\mathrm{e}^{3}+1).
\end{align*}
\end{solution}

\begin{example}\label{ex:6.4.3}
求 $\displaystyle\int_{0}^{1}\frac{\ln(1+x)}{1+x^{2}}\,\mathrm{d}x$.
\end{example}

\begin{solution}[解法 1]
\begin{align*}
\int_{0}^{1}\frac{\ln(1+x)}{1+x^{2}}\,\mathrm{d}x&\xlongequal{x=\tan t}\int_{0}^{\frac{\pi}{4}}\frac{\ln(1+\tan t)}{1+\tan^{2}t}\sec^{2}t\,\mathrm{d}t\\
&=\int_{0}^{\frac{\pi}{4}}\ln\frac{\cos t+\sin t}{\cos t}\,\mathrm{d}t=\int_{0}^{\frac{\pi}{4}}\ln\frac{\sqrt{2}\cos\left(\frac{\pi}{4}-t\right)}{\cos t}\,\mathrm{d}t\\
&=\int_{0}^{\frac{\pi}{4}}\ln\sqrt{2}\,\mathrm{d}t+\int_{0}^{\frac{\pi}{4}}\ln\cos\left(\frac{\pi}{4}-t\right)\mathrm{d}t-\int_{0}^{\frac{\pi}{4}}\ln\cos t\,\mathrm{d}t\\
&\xlongequal{u=\frac{\pi}{4}-t}\frac{\pi}{8}\ln2+\int_{\frac{\pi}{4}}^{0}\ln\cos u(-\mathrm{d}u)-\int_{0}^{\frac{\pi}{4}}\ln\cos t\,\mathrm{d}t\\
&=\frac{\pi}{8}\ln2+\int_{0}^{\frac{\pi}{4}}\ln\cos t\,\mathrm{d}t-\int_{0}^{\frac{\pi}{4}}\ln\cos t\,\mathrm{d}t=\frac{\pi}{8}\ln2.
\end{align*}
上述计算的妙处在于套用了换元的方法与步骤并利用定理~\ref{thm:6.4.1} 导致两个 Riemann 积分 $\int_{0}^{\frac{\pi}{4}}\ln\cos u\,\mathrm{d}u$ 与 $\int_{0}^{\frac{\pi}{4}}\ln\cos t\,\mathrm{d}t$ 相消, 从而得到 $\int_{0}^{1}\frac{\ln(1+x)}{1+x^{2}}\,\mathrm{d}x=\frac{\pi}{8}\ln2$. 但是, 无法用同样的换元来求出 $\int\frac{\ln(1+x)}{1+x^{2}}\,\mathrm{d}x$, 当然也无法用 Newton-Leibniz 公式.
\end{solution}

\begin{solution}[解法 2]
令
\[
I(\alpha)=\int_{0}^{1}\frac{\ln(1+\alpha x)}{1+x^{2}}\,\mathrm{d}x,
\]
其中 $x$ 为积分变量, $\alpha$ 为参变量. 则
\begin{align*}
I'(\alpha)&=\int_{0}^{1}\left(\frac{\ln(1+\alpha x)}{1+x^{2}}\right)'_{\alpha}\mathrm{d}x=\int_{0}^{1}\frac{x}{(1+\alpha x)(1+x^{2})}\,\mathrm{d}x\\
&=\frac{1}{1+\alpha^{2}}\int_{0}^{1}\left(\frac{-\alpha}{1+\alpha x}+\frac{x+\alpha}{1+x^{2}}\right)\mathrm{d}x\\
&=\frac{1}{1+\alpha^{2}}\left[-\ln(1+\alpha x)+\frac{1}{2}\ln(1+x^{2})+\alpha\arctan x\right]\Bigg|_{0}^{1}\\
&=-\frac{\ln(1+\alpha)}{1+\alpha^{2}}+\frac{\ln2}{2(1+\alpha^{2})}+\frac{\pi\alpha}{4(1+\alpha^{2})}.
\end{align*}
因此,
\begin{align*}
\int_{0}^{1}\frac{\ln(1+x)}{1+x^{2}}\,\mathrm{d}x&=I(1)=I(1)-I(0)\\
&\xlongequal{\text{Newton-Leibniz 公式}}\int_{0}^{1}I'(\alpha)\,\mathrm{d}\alpha\\
&=-\int_{0}^{1}\frac{\ln(1+\alpha)}{1+\alpha^{2}}\,\mathrm{d}\alpha+\frac{\ln2}{2}\int_{0}^{1}\frac{\mathrm{d}\alpha}{1+\alpha^{2}}+\frac{\pi}{4}\int_{0}^{1}\frac{\alpha}{1+\alpha^{2}}\,\mathrm{d}\alpha\\
&=-\int_{0}^{1}\frac{\ln(1+x)}{1+x^{2}}\,\mathrm{d}x+\frac{\ln2}{2}\arctan\alpha\Bigg|_{0}^{1}+\frac{\pi}{8}\ln(1+\alpha^{2})\Bigg|_{0}^{1}\\
&=-\int_{0}^{1}\frac{\ln(1+x)}{1+x^{2}}\,\mathrm{d}x+\frac{\pi\ln2}{8}+\frac{\pi}{8}\ln2,
\end{align*}
移项化简后得到
\[
\int_{0}^{1}\frac{\ln(1+x)}{1+x^{2}}\,\mathrm{d}x=\frac{1}{2}\cdot\frac{\pi}{4}\ln2=\frac{\pi}{8}\ln2.
\]
这个方法的妙处在于等号右边出现了 $-\int_{0}^{1}\frac{\ln(1+x)}{1+x^{2}}\,\mathrm{d}x$, 从而可解出 $\int_{0}^{1}\frac{\ln(1+x)}{1+x^{2}}\,\mathrm{d}x$.

但是, 遗憾的是
\[
I'(\alpha)=\int_{0}^{1}\left(\frac{\ln(1+\alpha x)}{1+x^{2}}\right)'_{\alpha}\mathrm{d}x,
\]
即求导与积分号可变换, 这里没有足够的理由, 读者学会了参变量积分求导的定理后再来检验它, 这等式确实是成立的.

我们给出解法 2, 是为了告诉读者除了解法 1 外还有一个很巧妙的方法, 还想告诉读者要思路开阔、大胆设想、不要受到束缚, 当然也不要忘了严格论证.
\end{solution}

\begin{example}\label{ex:6.4.4}
求 $\displaystyle\int_{0}^{\pi}\frac{x\sin x}{1+\cos^{2}x}\,\mathrm{d}x$.
\end{example}

\begin{solution}
\begin{align*}
\int_{0}^{\pi}\frac{x\sin x}{1+\cos^{2}x}\,\mathrm{d}x&\xlongequal{x=\pi-t}\int_{\pi}^{0}\frac{(\pi-t)\sin(\pi-t)}{1+\cos^{2}(\pi-t)}\,\mathrm{d}(\pi-t)\\
&=\int_{0}^{\pi}\frac{\pi\sin t}{1+\cos^{2}t}\,\mathrm{d}t+\int_{0}^{\pi}\frac{-t\sin t}{1+\cos^{2}t}\,\mathrm{d}t,
\end{align*}
移项化简得到
\begin{align*}
\int_{0}^{\pi}\frac{x\sin x}{1+\cos^{2}x}\,\mathrm{d}x&=\frac{\pi}{2}\int_{0}^{\pi}\frac{\sin x}{1+\cos^{2}x}\,\mathrm{d}x=-\frac{\pi}{2}\int_{0}^{\pi}\frac{\mathrm{d}\cos x}{1+\cos^{2}x}\\
&=-\frac{\pi}{2}\arctan(\cos x)\Bigg|_{0}^{\pi}=-\frac{\pi}{2}\left(\arctan(\cos\pi)-\arctan(\cos0)\right)\\
&=-\frac{\pi}{2}\left(-\frac{\pi}{4}-\frac{\pi}{4}\right)=\frac{\pi^{2}}{4}.
\end{align*}
\end{solution}

\begin{example}\label{ex:6.4.5}
证明:

(1) Wallis 公式
\[
\int_{0}^{\frac{\pi}{2}}\sin^{n}x\,\mathrm{d}x=\int_{0}^{\frac{\pi}{2}}\cos^{n}x\,\mathrm{d}x=\begin{cases}\dfrac{(n-1)!!}{n!!}\dfrac{\pi}{2}, & n\ \text{为偶数},\\[3mm]\dfrac{(n-1)!!}{n!!}, & n\ \text{为奇数};\end{cases}
\]
(2)
\[
\frac{\pi}{2}=\lim_{m\to+\infty}\left[\frac{(2m)!!}{(2m-1)!!}\right]^{2}\cdot\frac{1}{2m+1},
\]
其中
\[
(2m)!!=(2m)\cdot(2m-2)\cdots2,\quad(2m-1)!!=(2m-1)(2m-3)\cdots3\cdot1.
\]
\end{example}

\begin{proof}
(1)
\begin{align*}
I_{n}&=\int_{0}^{\frac{\pi}{2}}\sin^{n}x\,\mathrm{d}x=-\int_{0}^{\frac{\pi}{2}}\sin^{n-1}x\,\mathrm{d}\cos x\\
&\xlongequal{\text{分部积分}}-\sin^{n-1}x\cos x\Bigg|_{0}^{\frac{\pi}{2}}+\int_{0}^{\frac{\pi}{2}}(n-1)\sin^{n-2}x\cos^{2}x\,\mathrm{d}x\\
&=(n-1)\int_{0}^{\frac{\pi}{2}}\sin^{n-2}x(1-\sin^{2}x)\,\mathrm{d}x\\
&=(n-1)I_{n-2}-(n-1)I_{n},
\end{align*}
移项整理后得到递推公式
\[
I_{n}=\frac{n-1}{n}I_{n-2}\quad(n\geqslant2).
\]
由于
\[
I_{0}=\int_{0}^{\frac{\pi}{2}}\sin^{0}x\,\mathrm{d}x=\int_{0}^{\frac{\pi}{2}}\mathrm{d}x=\frac{\pi}{2},
\]
\[
I_{1}=\int_{0}^{\frac{\pi}{2}}\sin x\,\mathrm{d}x=-\cos x\Bigg|_{0}^{\frac{\pi}{2}}=\cos0-\cos\frac{\pi}{2}=1-0=1.
\]
重复应用递推公式便得
\begin{align*}
I_{2m}&=\frac{2m-1}{2m}I_{2m-2}=\frac{2m-1}{2m}\cdot\frac{2m-3}{2m-2}I_{2m-4}\\
&=\cdots\\
&=\frac{2m-1}{2m}\cdot\frac{2m-3}{2m-2}\cdots\frac{1}{2}I_{0}\\
&=\frac{(2m-1)!!}{(2m)!!}\frac{\pi}{2},
\end{align*}
\begin{align*}
I_{2m-1}&=\frac{2m-2}{2m-1}I_{2m-3}=\frac{2m-2}{2m-1}\cdot\frac{2m-4}{2m-3}I_{2m-5}\\
&=\cdots\\
&=\frac{2m-2}{2m-1}\cdot\frac{2m-4}{2m-3}\cdots\frac{2}{3}I_{1}\\
&=\frac{(2m-2)!!}{(2m-1)!!}.
\end{align*}
而
\begin{align*}
\int_{0}^{\frac{\pi}{2}}\cos^{n}x\,\mathrm{d}x&\xlongequal{x=\frac{\pi}{2}-t}-\int_{\frac{\pi}{2}}^{0}\cos^{n}\left(\frac{\pi}{2}-t\right)\mathrm{d}t\\
&=\int_{0}^{\frac{\pi}{2}}\sin^{n}t\,\mathrm{d}t=\int_{0}^{\frac{\pi}{2}}\sin^{n}x\,\mathrm{d}x\\
&=\begin{cases}\dfrac{(n-1)!!}{n!!}\cdot\dfrac{\pi}{2}, & n\ \text{为偶数},\\[
3mm]\dfrac{(n-1)!!}{n!!}, & n\ \text{为奇数}.\end{cases}
\end{align*}
或类似 $\int_{0}^{\frac{\pi}{2}}\sin^{n}x\,\mathrm{d}x$, 应用分部积分法得到递推公式求得上述公式.
(2) 因为
\[
\int_{0}^{\frac{\pi}{2}}\sin^{2m+1}x\,\mathrm{d}x<\int_{0}^{\frac{\pi}{2}}\sin^{2m}x\,\mathrm{d}x<\int_{0}^{\frac{\pi}{2}}\sin^{2m-1}x\,\mathrm{d}x.
\]
所以由 (1), 上式就是
\[
\frac{(2m)!!}{(2m+1)!!}<\frac{(2m-1)!!}{(2m)!!}\cdot\frac{\pi}{2}<\frac{(2m-2)!!}{(2m-1)!!}.
\]
由此得到
\[
A_{m}=\frac{1}{2m+1}\left[\frac{(2m)!!}{(2m-1)!!}\right]^{2}<\frac{\pi}{2}<\left[\frac{(2m)!!}{(2m-1)!!}\right]^{2}\cdot\frac{1}{2m}=B_{m}.
\]
因为
\begin{align*}
0<B_{m}-A_{m}&=\left[\frac{(2m)!!}{(2m-1)!!}\right]^{2}\left(\frac{1}{2m}-\frac{1}{2m+1}\right)\\
&=\left[\frac{(2m)!!}{(2m-1)!!}\right]^{2}\frac{1}{2m(2m+1)}<\frac{1}{2m}\cdot\frac{\pi}{2}\to0\quad(m\to+\infty),
\end{align*}
故由夹逼定理立即有
\[
\lim_{m\to+\infty}(B_{m}-A_{m})=0.
\]
再由 $0<\frac{\pi}{2}-A_{m}<B_{m}-A_{m}\to0$ $(m\to+\infty)$ 得到
\[
\lim_{m\to+\infty}\left(\frac{\pi}{2}-A_{m}\right)=0,
\]
\begin{align*}
\lim_{m\to+\infty}\left[\frac{(2m)!!}{(2m-1)!!}\right]^{2}\frac{1}{2m+1}&=\lim_{m\to+\infty}A_{m}=\lim_{m\to+\infty}\left[\frac{\pi}{2}-\left(\frac{\pi}{2}-A_{m}\right)\right]\\
&=\frac{\pi}{2}-0=\frac{\pi}{2}.
\end{align*}
\end{proof}

\begin{example}\label{ex:6.4.6}
计算 $\displaystyle I(m,n)=\int_{0}^{\frac{\pi}{2}}\cos^{m}x\sin^{n}x\,\mathrm{d}x$.
\end{example}

\begin{solution}
\begin{align*}
I(m,n)&=\frac{1}{n+1}\int_{0}^{\frac{\pi}{2}}\cos^{m-1}x\,\mathrm{d}\sin^{n+1}x\\
&=\frac{1}{n+1}\left[\cos^{m-1}x\sin^{n+1}x\Bigg|_{0}^{\frac{\pi}{2}}-\int_{0}^{\frac{\pi}{2}}(m-1)\cos^{m-2}x(-\sin^{n+2}x)\,\mathrm{d}x\right]\\
&=\frac{m-1}{n+1}\int_{0}^{\frac{\pi}{2}}\cos^{m-2}x\sin^{n}x(1-\cos^{2}x)\,\mathrm{d}x\\
&=\frac{m-1}{n+1}\left[I(m-2,n)-I(m,n)\right],
\end{align*}
移项化简得
\begin{align*}
I(m,n)&=\frac{m-1}{m+n}I(m-2,n)\\
&=\begin{cases}\dfrac{(m-1)(m-3)\cdots1}{(m+n)(m+n-2)\cdots(n+2)}I(0,n), & m\ \text{为偶数},\\[
3mm]\dfrac{(m-1)(m-3)\cdots2}{(m+n)(m+n-2)\cdots(n+3)}I(1,n), & m\ \text{为奇数},\end{cases}\\
&=\begin{cases}\dfrac{(m-1)!!(n-1)!!}{(m+n)!!}\dfrac{\pi}{2}, & m\ \text{与}\ n\ \text{均为偶数},\\[3mm]\dfrac{(m-1)!!(n-1)!!}{(m+n)!!}, & m\ \text{与}\ n\ \text{中有一个奇数}.\end{cases}
\end{align*}
\end{solution}
\begin{example}\label{ex:6.4.7}
求 $\displaystyle\lim_{n\to+\infty}\left[\left(1+\frac{1}{n}\right)\left(1+\frac{2}{n}\right)\cdots\left(1+\frac{n}{n}\right)\right]^{\frac{1}{n}}$.
\end{example}
\begin{solution}
因为
\begin{align*}
\lim_{n\to+\infty}\frac{1}{n}\sum_{k=1}^{n}\left[\ln\left(1+\frac{k}{n}\right)\right]&=\int_{0}^{1}\ln(1+x)\,\mathrm{d}x=\int_{0}^{1}\ln(1+x)\,\mathrm{d}(1+x)\\
&\xlongequal{\text{分部积分}}(1+x)\ln(1+x)\Bigg|_{0}^{1}-\int_{0}^{1}(1+x)\cdot\frac{1}{1+x}\,\mathrm{d}x\\
&=2\ln2-\int_{0}^{1}\mathrm{d}x=2\ln2-1,
\end{align*}
所以
\begin{align*}
&\lim_{n\to+\infty}\left[\left(1+\frac{1}{n}\right)\left(1+\frac{2}{n}\right)+\cdots+\left(1+\frac{n}{n}\right)\right]^{\frac{1}{n}}\\
&=\exp\left[\lim_{n\to+\infty}\frac{1}{n}\sum_{k=1}^{n}\ln\left(1+\frac{k}{n}\right)\right]=\mathrm{e}^{2\ln2-1}=\frac{\mathrm{e}^{\ln4}}{\mathrm{e}}=\frac{4}{\mathrm{e}}.
\end{align*}
\end{solution}
\begin{example}\label{ex:6.4.8}
计算 $\displaystyle\int_{0}^{\frac{\pi}{4}}\frac{\mathrm{d}x}{\cos x}$.
\end{example}
\begin{solution}
\begin{align*}
\int_{0}^{\frac{\pi}{4}}\frac{\mathrm{d}x}{\cos x}&=\int_{0}^{\frac{\pi}{4}}\frac{\mathrm{d}\sin x}{1-\sin^{2}x}\xlongequal{t=\sin x}\int_{0}^{\frac{\sqrt{2}}{2}}\frac{\mathrm{d}t}{1-t^{2}}\\
&=\frac{1}{2}\int_{0}^{\frac{\sqrt{2}}{2}}\left(\frac{1}{1+t}+\frac{1}{1-t}\right)\mathrm{d}t\\
&=\frac{1}{2}\ln\frac{1+t}{1-t}\Bigg|_{0}^{\frac{\sqrt{2}}{2}}=\frac{1}{2}\ln\frac{1+\frac{\sqrt{2}}{2}}{1-\frac{\sqrt{2}}{2}}=\ln(\sqrt{2}+1).
\end{align*}
\end{solution}
\begin{example}\label{ex:6.4.9}
(1) 设 $f$ 为连续的奇函数, 则对 $\forall a\in\mathbb{R}$, 有
\[
\int_{-a}^{a}f(x)\,\mathrm{d}x=0;
\]
(2) 设 $f$ 为连续的偶函数, 则对 $\forall a\in\mathbb{R}$, 有
\[
\int_{-a}^{a}f(x)\,\mathrm{d}x=2\int_{0}^{a}f(x)\,\mathrm{d}x;
\]
(3) 设 $f$ 是周期为 $T$ 的连续函数, 则对 $\forall a\in\mathbb{R}$, 有
\[
\int_{a}^{a+T}f(x)\,\mathrm{d}x=\int_{0}^{T}f(x)\,\mathrm{d}x.
\]
\end{example}

\begin{proof}
(1)
\begin{align*}
\int_{-a}^{a}f(x)\,\mathrm{d}x&=\int_{-a}^{0}f(x)\,\mathrm{d}x+\int_{0}^{a}f(x)\,\mathrm{d}x\xlongequal{x=-t}\int_{a}^{0}f(-t)\,\mathrm{d}(-t)+\int_{0}^{a}f(x)\,\mathrm{d}x\\
&=\int_{0}^{a}-f(t)\,\mathrm{d}t+\int_{0}^{a}f(x)\,\mathrm{d}x=0.
\end{align*}
(2)
\begin{align*}
\int_{-a}^{a}f(x)\,\mathrm{d}x&=\int_{-a}^{0}f(x)\,\mathrm{d}x+\int_{0}^{a}f(x)\,\mathrm{d}x\\
&\xlongequal{x=-t}\int_{a}^{0}f(-t)\,\mathrm{d}(-t)+\int_{0}^{a}f(x)\,\mathrm{d}x\\
&=\int_{0}^{a}f(t)\,\mathrm{d}t+\int_{0}^{a}f(x)\,\mathrm{d}x=2\int_{0}^{a}f(x)\,\mathrm{d}x.
\end{align*}
(3) 因为
\begin{align*}
\int_{T}^{a+T}f(x)\,\mathrm{d}x&\xlongequal{x=t+T}\int_{0}^{a}f(t+T)\,\mathrm{d}(t+T)=\int_{0}^{a}f(t)\,\mathrm{d}t\\
&=\int_{0}^{a}f(x)\,\mathrm{d}x,
\end{align*}
所以
\begin{align*}
\int_{a}^{a+T}f(x)\,\mathrm{d}x&=\int_{a}^{0}f(x)\,\mathrm{d}x+\int_{0}^{T}f(x)\,\mathrm{d}x+\int_{T}^{a+T}f(x)\,\mathrm{d}x\\
&=-\int_{0}^{a}f(x)\,\mathrm{d}x+\int_{0}^{T}f(x)\,\mathrm{d}x+\int_{0}^{a}f(x)\,\mathrm{d}x\\
&=\int_{0}^{T}f(x)\,\mathrm{d}x.
\end{align*}
\end{proof}

\begin{remark}\label{rem:6.4.1}
在例~\ref{ex:6.4.9} 中, 如果将 $f$「连续」改为「Riemann 可积」, 则其结论仍正确. 只需应用 Riemann 积分定义证明:
\[
(1)\ \int_{0}^{a}f(x)\,\mathrm{d}x=-\int_{-a}^{0}f(x)\,\mathrm{d}x;\qquad(2)\ \int_{0}^{a}f(x)\,\mathrm{d}x=\int_{-a}^{0}f(x)\,\mathrm{d}x;
\]
\[
(3)\ \int_{T}^{a+T}f(x)\,\mathrm{d}x=\int_{0}^{a}f(x)\,\mathrm{d}x.
\]
或者应用定理~\ref{thm:6.4.1pp} 及例~\ref{ex:6.4.9} 中的方法.
\end{remark}

\begin{exblock}{练习题 6.4}

  \item 证明: $\lim\limits_{n\to+\infty}\prod\limits_{k=1}^{n}\left(1+\dfrac{k}{n}\right)^{\frac{k}{n^{2}}}=\mathrm{e}^{\frac{1}{4}}$.

  (提示: $\lim\limits_{n\to+\infty}\dfrac{1}{n}\sum\limits_{k=1}^{n}\dfrac{k}{n}\ln\left(1+\dfrac{k}{n}\right)=\int_{0}^{1}x\ln(1+x)\,\mathrm{d}x$.)

  \item 求下列极限:
  \begin{enumerate}[label=(\arabic*)]
    \item $\lim\limits_{n\to+\infty}\sin\dfrac{\pi}{n}\sum\limits_{k=1}^{n}\left(2+\cos\dfrac{k\pi}{n}\right)^{-1}$;
    \item $\lim\limits_{n\to+\infty}\sum\limits_{i=1}^{n}\dfrac{i}{n^{2}}\sin\left(\dfrac{i\pi}{2n}\right)^{2}$.
  \end{enumerate}

  \item 求下列 Riemann 积分:
  \begin{enumerate}[label=(\arabic*)]
    \item $\int_{0}^{1}x(2-x^{2})^{\frac{1}{2}}\,\mathrm{d}x$;
    \item $\int_{0}^{\frac{\pi}{2}}\dfrac{\mathrm{d}x}{1+\varepsilon\cos x}$.
  \end{enumerate}

  \item 证明: $\int_{\mathrm{e}^{-2n\pi}}^{1}\left|\left[\cos\left(\ln\dfrac{1}{x}\right)\right]'\right|\mathrm{d}x=4n$, $n\in\mathbb{N}$.

  \item 设 $I$ 为一个开区间, 函数 $f$ 在 $I$ 内连续, $a<b$ 且 $a,b\in I$. 证明:
  \[
  \lim_{h\to0}\frac{1}{h}\int_{a}^{b}[f(x+h)-f(x)]\,\mathrm{d}x=f(b)-f(a).
\]
  (提示: 先作换元 $t=x+h$, 再应用积分中值定理.)

  \item 计算下列 Riemann 积分:
  \begin{enumerate}[label=(\arabic*)]
    \item $\int_{0}^{\pi}\sin^{3}x\,\mathrm{d}x$;
    \item $\int_{-\pi}^{\pi}x^{2}\cos x\,\mathrm{d}x$;
    \item $\int_{0}^{\sqrt{3}}x\arctan x\,\mathrm{d}x$;
    \item $\int_{0}^{3}\dfrac{x}{1+\sqrt{1+x}}\,\mathrm{d}x$;
    \item $\int_{-1}^{0}(2x+1)\sqrt{1-x-x^{2}}\,\mathrm{d}x$;
    \item $\int_{\mathrm{e}^{-1}}^{\mathrm{e}}|\ln x|\,\mathrm{d}x$;
    \item $\int_{0}^{5}[x]\sin\dfrac{\pi x}{5}\,\mathrm{d}x$;
    \item $\int_{0}^{a}x^{2}\sqrt{a^{2}-x^{2}}\,\mathrm{d}x$;
    \item $\int_{0}^{\ln2}\sqrt{\mathrm{e}^{x}-1}\,\mathrm{d}x$;
    \item $\int_{0}^{1}x^{n}\ln x\,\mathrm{d}x$, $n\in\mathbb{N}$;
    \item $\int_{0}^{\frac{\pi}{2}}\dfrac{\cos x\sin x}{a^{2}\sin^{2}x+b^{2}\cos^{2}x}\,\mathrm{d}x$, $ab\neq0$;
    \item $\int_{0}^{a}\ln\left(x+\sqrt{a^{2}+x^{2}}\right)\mathrm{d}x$, $a>0$.
  \end{enumerate}

  \item 证明: $\lim\limits_{n\to+\infty}\int_{0}^{\frac{\pi}{2}}\sin^{n}x\,\mathrm{d}x=0$. 由此推证
  \begin{enumerate}[label=(\arabic*)]
    \item $\lim\limits_{n\to+\infty}\dfrac{(2n)!!}{(2n+1)!!}=0$;
    \item $\lim\limits_{n\to+\infty}\dfrac{(2n-1)!!}{(2n)!!}=0$.
  \end{enumerate}

  \item 设 $f$ 为连续函数, 证明:
  \begin{enumerate}[label=(\arabic*)]
    \item $\int_{0}^{\frac{\pi}{2}}f(\cos x)\,\mathrm{d}x=\int_{0}^{\frac{\pi}{2}}f(\sin x)\,\mathrm{d}x$;
    \item $\int_{0}^{\pi}xf(\sin x)\,\mathrm{d}x=\dfrac{\pi}{2}\int_{0}^{\pi}f(\sin x)\,\mathrm{d}x$.
  \end{enumerate}
  并由 (2) 证明 $\int_{0}^{\pi}\dfrac{x\sin x}{1+\cos^{2}x}\,\mathrm{d}x=\dfrac{\pi^{2}}{4}$.

  \item 计算下列 Riemann 积分:
  \begin{enumerate}[label=(\arabic*)]
    \item $\int_{0}^{\frac{\pi}{2}}\dfrac{\cos^{2}x}{\cos x+\sin x}\,\mathrm{d}x$;
    \item $\int_{0}^{\frac{\pi}{2}}\dfrac{\sin^{2}x}{\cos x+\sin x}\,\mathrm{d}x$.
  \end{enumerate}

  \item 设 $f$ 在 $[a,b]$ 上连续可导, 应用分部积分证明:
  \begin{enumerate}[label=(\arabic*)]
    \item $\lim\limits_{\lambda\to+\infty}\int_{a}^{b}f(x)\cos\lambda x\,\mathrm{d}x=0$;
    \item $\lim\limits_{\lambda\to+\infty}\int_{a}^{b}f(x)\sin\lambda x\,\mathrm{d}x=0$.
  \end{enumerate}
  若将 $f$ 放宽成「在 $[a,b]$ 上 Riemann 可积」, 则结论仍然正确. 那就是著名的 Riemann 引理. 本书第 3 册中将讨论它.

  \item 设 $f$ 在 $[0,1]$ 上 2 阶连续可导, 证明:
  \[
  \int_{0}^{1}f(x)\,\mathrm{d}x=\frac{1}{2}\left[f(0)+f(1)-\int_{0}^{1}x(1-x)f''(x)\,\mathrm{d}x\right].
\]

  \item 应用分部积分证明: $\int_{0}^{2\pi}\left(\int_{x}^{2\pi}\frac{\sin t}{t}\,\mathrm{d}t\right)\mathrm{d}x=0$.

  \item 设 $f$ 在 $[a,+\infty)$ 中连续, 对 $\forall a>0$, 证明:
  \[
  \int_{0}^{a}\left(\int_{0}^{x}f(t)\,\mathrm{d}t\right)\mathrm{d}x=\int_{0}^{a}f(x)(a-x)\,\mathrm{d}x.
\]

  \item 计算下列 Riemann 积分:
  \begin{enumerate}[label=(\arabic*)]
    \item $\int_{0}^{\frac{\pi}{2}}\cos^{5}x\sin2x\,\mathrm{d}x$;
    \item $\int_{0}^{1}\sqrt{4-x^{2}}\,\mathrm{d}x$;
    \item $\int_{0}^{a}x^{2}\sqrt{a^{2}-x^{2}}\,\mathrm{d}x$ $(a>0)$;
    \item $\int_{0}^{1}\dfrac{\mathrm{d}x}{(x^{2}-x+1)^{\frac{3}{2}}}$;
    \item $\int_{0}^{1}\dfrac{\mathrm{d}x}{\mathrm{e}^{x}+\mathrm{e}^{-x}}$;
    \item $\int_{0}^{\frac{\pi}{2}}\dfrac{\cos x}{1+\sin^{2}x}\,\mathrm{d}x$;
    \item $\int_{0}^{1}\arcsin x\,\mathrm{d}x$;
    \item $\int_{0}^{\frac{\pi}{2}}\mathrm{e}^{x}\sin x\,\mathrm{d}x$;
    \item $\int_{0}^{1}\mathrm{e}^{\sqrt{x}}\,\mathrm{d}x$;
    \item $\int_{0}^{a}x^{2}\sqrt{\dfrac{a-x}{a+x}}\,\mathrm{d}x$;
    \item $\int_{0}^{\frac{\pi}{2}}\dfrac{\cos\theta}{\sin\theta+\cos\theta}\,\mathrm{d}\theta$.
  \end{enumerate}

\end{exblock}

\begin{exblock}{思考题 6.4}

  \item 设 $m,n\in\mathbb{N}$, $B(m,n)=\sum\limits_{k=0}^{n}C_{n}^{k}\dfrac{(-1)^{k}}{m+k+1}$. 证明:
  \begin{enumerate}[label=(\arabic*)]
    \item $B(m,n)=\int_{0}^{1}x^{m}(1-x)^{n}\,\mathrm{d}x$;
    \item $B(m,n)=B(n,m)$;
    \item $B(m,n)=\dfrac{m!n!}{(m+n+1)!}$.
  \end{enumerate}

  \item 对于 $\alpha\in(0,1]$, 定义 $f_{\alpha}(x)=\left[\dfrac{\alpha}{x}\right]-\alpha\left[\dfrac{1}{x}\right]$. 证明:
  \[
  \int_{0}^{1}f_{\alpha}(x)\,\mathrm{d}x=\alpha\ln\alpha.
\]
  (提示: $f_{\alpha}(x)=-\left(\dfrac{\alpha}{x}-\left[\dfrac{\alpha}{x}\right]\right)+\alpha\left(\dfrac{1}{x}-\left[\dfrac{1}{x}\right]\right)$, 对积分 $\int_{0}^{\alpha}\left(\dfrac{\alpha}{x}-\left[\dfrac{\alpha}{x}\right]\right)\mathrm{d}x$ 作换元 $x=\alpha t$.)

  \item 设 $I(m,n)=\int_{0}^{\frac{\pi}{2}}\cos^{m}x\sin^{n}x\,\mathrm{d}x$ ($m,n\in\mathbb{N}$). 证明:
  \[
  I(m,n)=\frac{m-1}{m+n}I(m-2,n)=\frac{n-1}{m+n}I(m,n-2).
\]
  并求 $I(2m,2n)$.

  \item 设 $x>0$, 证明不等式:
  \[
  \left|\int_{x}^{x+c}\sin t^{2}\,\mathrm{d}t\right|<\frac{1}{x},\quad\text{其中}\ c>0.
\]

\end{exblock}
